\subsection{Sequencer Boundary Detection Scaling}
\label{sec:sequencer}

\paragraph{Task Description.}
The sequencer experiment uses boundary-aware splits (32-byte suffix
holdouts) on a 200-sample corpus. $N = 50$ queries are issued in
corpus-position order. Does boundary detection maintain prediction
quality as queries target progressively deeper data? Metrics are
beam confidence, rejection rate, and origin diversity.

\paragraph{Results.}
Beam confidence remains stable across 50 queries (top-score $$\mu = -0.003$, $\sigma = 0.020$, median $= -0.007$, range $[-0.033,\,0.032]$, $N = 50$$).
Rejection rate: $\mu = 0.875$. Origin diversity: $\mu = 1.0$ tries.
Substrate entries: 50.

Figure~\ref{fig:sequencer_map} shows beam confidence trajectory
(top-left), cumulative convergence (top-right), boundary routing
noise (bottom-left), and mesh utilization (bottom-right).

\paragraph{Assessment.}
High rejection ($0.875$) despite boundary-aware splits suggests the
sequencer's boundary detection is misaligned with the trie mesh's
attractor structure. The boundary-chosen split points produce
prefixes that don't match any trie's centroid well. The phase
concentration signal indicates whether this is correctable by
mesh re-balancing or fundamental to the byte distribution.
